% part: intuitionistic-logic
% chapter: propositions-as-types
% section: type-preservation

\documentclass[../../../include/open-logic-section]{subfiles}

\begin{document}

\olfileid{int}{pty}{tpr}

\olsection{حفظ النوع}

يمكن أيضًا تعريف تحويلات الاختزال والتبديل المدروسة مباشرةً للنسقين
\Log{tN2i} و\Log{tN3i}. ويعرّف النسق~\Log{tN3i} نوع حد لامبدا~$N$ بالنسبة
إلى السياق~$\Gamma$ بأنه~!!{formula}~$!A$ التي يكون للمتتالية
$\Gamma\Sequent N:!A$~!!a{proof} في~\Log{tN3i}. ولأن قواعد اختزال بيتا
لحدود لامبدا تطابق تمامًا تحويلات الاختزال والتبديل على~!!{proof}s، تنتج
نتيجة مهمة: يبقى نوع حد لامبدا بالنسبة إلى السياق~$\Gamma$ على حاله بعد
كل خطوة واحدة من اختزال بيتا. وتسمى هذه الخاصية \emph{حفظ النوع}.

\begin{prop}
إذا كان نوع~$N$ هو~$!A$ بالنسبة إلى~$\Gamma$، وكان~$N\redone N'$ في خطوة
اختزال واحدة، فإن نوع~$N'$ هو~$!A$ بالنسبة إلى~$\Gamma$.
\end{prop}

\begin{proof}
  نستطيع بسهولة أن نثبت، بالاستقراء في~$N$ وفي ارتفاع~!!{proof}s~$\delta$
  للمتتالية~$\Gamma\Sequent N:!A$، ما يأتي: إذا كان~$M$ حدًا جزئيًا من
  $N$، احتوى~$\delta$~!!{proof}s جزئيًا~$\delta_1$ ينتهي بـ
  $\Gamma'\Sequent M:!B$. وإذا كان~$M$ ردكسًا، انطبق تحويل على
  $\delta_1$. وبتطبيق الاختزال على~$\delta_1$ نحصل على~!!a{proof}
  $\delta_1'$ لـ~$\Gamma'\Sequent M':!B$. وبفحص تحويلات~\Log{tN3i} نجد
  أن~$M'$ هو مختزَل~$M$. وإذا استبدلنا~$\delta_1'$ بـ~$\delta_1$ في
  $\delta$ نتج~!!a{proof}~$\delta'$ لـ~$\Gamma\Sequent N':!A$، حيث ينتج
  $N'$ من~$N$ باستبدال~$M'$ بالحد الجزئي~$M$.

  انظر، مثلًا، في تحويل الاختزال لـ~\Intro{\land} يتبعه~\Elim{\land}:
  \[
  \bottomAlignProof
  \AxiomC{}
  \Deduce$\Gamma \fCenter N_1 : !B_1$
  \AxiomC{}
  \Deduce$\Gamma \fCenter N_2 : !B_2$
  \RightLabel{\Intro{\land}}
  \BinaryInf$\Gamma \fCenter \pair{N_1}{N_2} : !B_1 \land !B_2$
  \RightLabel{\Elim{\land}[i]}
  \UnaryInf$\Gamma \fCenter \proj{i}{\pair{N_1}{N_2}} : !B_i$
  \DisplayProof
  \quad\redone\quad
  \bottomAlignProof
  \AxiomC{}
  \Deduce$\Gamma \fCenter N_i : !B_i$
  \DisplayProof
  \]
  نرى أن حد برهان~!!{proof} الواقع على اليمين، أي~$\delta_1'$، هو~$N_i$.
  وهو بالضبط نتيجة تطبيق خطوة اختزال بيتا على حد البرهان
  $\proj{i}{\pair{N_1}{N_2}}$ لـ~!!{proof} الواقع على اليسار، أي
  $\delta_1$.

  أو انظر في استدلال~\Intro{\lif} يتبعه استدلال~\Elim{\lif}، وفي تحويل
  الاختزال المطبق عليه:
  \begin{multline*}
  \bottomAlignProof
  \AxiomC{$\delta_2$}
  \Deduce$x: !B_1, \Gamma \fCenter N : !B_2$
  \RightLabel{\Intro{\lif}}
  \UnaryInf$\Gamma \fCenter \lambd[\typeof{x}{!B_1}][N] : !B_1 \lif !B_2$
  \AxiomC{}
  \RightLabel{$\delta_3$}
  \Deduce$\Gamma \fCenter M : !B_1$
  \RightLabel{\Elim{\lif}}
  \BinaryInf$\Gamma \fCenter (\lambd[\typeof{x}{!B_1}][N] M) : !B_2$
  \DisplayProof
  \quad\redone\\
  \bottomAlignProof
  \AxiomC{}
  \RightLabel{$\delta_3$}
  \Deduce$\Gamma \fCenter M : !B_1$
  \RightLabel{$\Subst{\delta_2}{\delta_3}{x:!B_1}$}
  \Deduce$\Gamma \fCenter \Subst{N}{M}{x} : !B_2$
  \DisplayProof
  \end{multline*}
  ومرة أخرى يكون حد برهان~!!{proof} الواقع على اليمين هو بالضبط نتيجة
  اختزال حد برهان~!!{proof} الواقع على اليسار ببيتا خطوةً واحدة. ويبقى أن
  نتحقق بعناية، بالاستقراء في ارتفاع~$\delta_2$، من أن استبدال~!!{proof}
  $\delta_3$ لـ~$\Gamma\Sequent M:!B_1$ بكل بديهية من الصورة
  $\Gamma'\Sequent x:!B_1$ في~$\delta_2$ ينتج فعلًا~!!a{proof}
  $\Subst{\delta_2}{\delta_3}{x:!B_1}$ للمتتالية
  $\Gamma\Sequent\Subst{N}{M}{x}:!B_2$.
\end{proof}

وللتطابق الوثيق بين~!!{proof}s وحدود لامبدا نتيجة أخرى. فإذا كان~$N$ حد
لامبدا \emph{مغلقًا} من النوع~$!A$ وكان له ردكس~$M$، احتوى البرهان
$\delta$ المقابل في~\Log{tN3i} للمتتالية~$\Sequent N:!A$~!!{proof} جزئيًا
ينتهي بـ~$\Gamma'\Sequent M:!B$ وينطبق عليه تحويل. وإذا كان
$N\redone N'$ باستبدال مختزَل~$M$ به، كانت نتيجة تطبيق الاختزال المقابل
على~$\delta$~!!a{proof} لـ~$\Sequent N':!A$.

أما خاصية \emph{التقدم} الصحيحة فتخص الحدود المغلقة الصحيحة النوع: يبيّن
الاستقراء في اشتقاق التنميط، باستعمال لمّة الصيغ القانونية للقيم، أن كل حد
مغلق صحيح النوع إما قيمة ذات الصيغة القانونية لنوعه، وإما يختزل في خطوة
واحدة إلى حد آخر. ويضمن حفظ النوع أن الحد الناتج له النوع نفسه. ولذلك لا
``تعلق'' عملية اختزال بيتا لحدود مغلقة صحيحة النوع؛ أما الحدود المفتوحة فقد
تكون في صيغة عادية محايدة، فلا يشملها ادعاء التقدم هذا.

\end{document}
